% Generated draft: requires mathematical and source review.
\documentclass[11pt]{article}
\usepackage[margin=25mm]{geometry}
\usepackage{fontspec}
\setmainfont{Latin Modern Roman}
\usepackage{amsmath,amssymb,mathtools,unicode-math}
\setmathfont{Latin Modern Math}
\usepackage{enumitem,needspace,xurl,hyperref}
\setlength{\parindent}{0pt}
\setlength{\parskip}{5pt}
\setlength{\emergencystretch}{3em}
\newcommand{\N}{\mathbb{N}}
\newcommand{\Z}{\mathbb{Z}}
\newcommand{\Q}{\mathbb{Q}}
\newcommand{\R}{\mathbb{R}}
\newcommand{\C}{\mathbb{C}}
\newcommand{\F}{\mathbb{F}}
\newcommand{\A}{\mathbb{A}}
\newcommand{\PP}{\mathbb{P}}
\newcommand{\E}{\mathbb{E}}
\newcommand{\eps}{\varepsilon}
\newcommand{\dd}{\,\mathrm{d}}
\begin{document}
\section*{UnsolvedMath ID: 800656}
\textbf{Source problem code:} catalogue entry
\subsection{Definitions and mathematical statement}
Let \(p\) be an ultrafilter on the sequence index set \(\N\). For a sequence \((y_n)_{n\in\N}\) in a topological space \(Y\), define
\[
y=p\text{-}\lim_{n\in\N}y_n
\quad\Longleftrightarrow\quad
\{n\in\N:y_n\in U\}\in p
\text{ for every neighborhood }U\text{ of }y.
\]
Using the usual characterization of sequential spaces by sequentially closed subsets, the supplied definition of \(p\)-sequentiality reads as follows: \(Y\) is \(p\)-sequential if every subset \(A\subseteq Y\) satisfying
\[
\forall (a_n)_{n\in\N}\in A^{\N}\ \forall y\in Y,
\qquad y=p\text{-}\lim_{n\in\N}a_n\ \Longrightarrow\ y\in A
\]
is closed in \(Y\).

Under the conventional real-valued interpretation, \(C_p(X)\) denotes the set \(C(X,\R)\) of continuous functions with the topology inherited from the product \(\R^X\). The subscript in \(C_p(X)\) denotes pointwise convergence, whereas the ultrafilter \(p\) is the parameter in \(p\)-sequentiality. The supplied record does not specify the separation assumptions on \(X\) or explicitly confirm this function-space convention.

An open \(\omega\)-cover of \(X\), under the usual convention, is a family \(\mathcal U\) of open subsets of \(X\) such that \(X\notin\mathcal U\) and every finite subset of \(X\) is contained in some member of \(\mathcal U\). Write \(\Omega(X)\) for the intended class of such covers; whether this class is restricted to countable covers is not stated in the record. For classes of covers \(\mathscr A\) and \(\mathscr B\), the usual interpretation of \(\binom{\mathscr A}{\mathscr B}\) is that every member of \(\mathscr A\) has a subfamily belonging to \(\mathscr B\).

The supplied question asks whether, for the intended spaces \(X\) and ultrafilters \(p\),
\[
C_p(X)\text{ is }p\text{-sequential}
\quad\Longrightarrow\quad
\forall\mathcal U\in\Omega(X)\ \exists\mathcal V\subseteq\mathcal U,
\qquad \mathcal V\in\mathcal O_p(X).
\]
Here \(\mathcal O_p(X)\) retains the source notation: its definition is absent from the supplied record. The record also does not say whether \(p\) must be free or whether the implication is intended for every ultrafilter in a particular class. Consequently, this is a statement schema with unresolved source qualifications, rather than a fully specified assertion.
\subsection{Short English statement}
The question asks whether \(p\)-sequentiality of the pointwise function space \(C_p(X)\) forces every open \(\omega\)-cover of \(X\) to contain a subfamily in \(\mathcal O_p(X)\). The supplied record leaves \(\mathcal O_p\), the admissible spaces and ultrafilters, and the cover-countability convention unspecified.
\subsection{Sources}
\begin{itemize}
\item Local UnsolvedMath registry, permanent ID \(800656\), exact source problem code \texttt{catalogue entry}; supplied complete source record. The registry attributes the problem to Tsaban, with proposed year 2003, and labels it open. This status is supplied metadata, not an independently verified mathematical assessment.
\item Tsaban, \emph{Selection Principles in Mathematics: A Milestone of Open Problems} (2003), Problem 29. Supplied source URL: \url{https://arxiv.org/abs/math/0312182}. The record states extraction method \texttt{arxiv\_latex} and access date 2026-10-16; these are supplied provenance claims, not a claim of access during preparation of this draft.
\item Supplied attribution/context labels: [GT1] and [GTpseq, Problem 0.10]. The record supplies no bibliographic expansions or separate URLs for these references.
\end{itemize}
\end{document}
